%% GENERATED by tools/from_docs.py from stepss-docs. Do not edit by hand:
%% edit the page named below in stepss-docs and regenerate.
%% source: stepss-docs/src/content/docs/models/governors/plant.mdx

\chapter{Plant turbine-governor models}\label{doc:model-tor-plant}

This page documents \texttt{GASTURBM}, \texttt{GOVCLASM}, \texttt{GOVHYDR} and \texttt{GOVNUC}, which the family dispatcher has handled directly since 3.79. Before 3.79 they were compiled under no name and RAMSES rejected them in a data file. They model a gas turbine, a classical steam plant, a hydro plant and a nuclear plant.

The \texttt{TOR} sub-record these models are entered with is described on the \hyperref[doc:model-tor]{Governor Models} page.

\section{tor\_gasturbm}\label{model-tor-plant:tor_gasturbm}

\begin{docsnote}{Note}

Callable since 3.79. Earlier releases compiled this model under no name, so a data file referring to it was rejected.

\end{docsnote}

\subsection{Scientific Description}\label{model-tor-plant:scientific-description}

The \texttt{tor\_gasturbm} model is a \textbf{gas turbine} governor following a multi-stage combustion model. It captures the compressor, fuel system, combustion chamber delay, and turbine torque characteristic. The name suffix \texttt{m} indicates a modified version.

\textbf{Speed governor:}

The governor error drives a PI controller producing a fuel demand signal. The speed governor uses:

\[e_{gov} = (\omega_{ref} - \omega) / TDSP - STAT \cdot P_{elec}\]

with saturation limits \([MIN,\, MAX]\).

\textbf{Fuel system and combustion:}

The fuel valve position \texttt{g1} is driven through:
- A speed governor with gain \(ZZ\) and lead--lag time constants determined by \(XX\), \(YY\)
- A fuel flow path through first-order lag \(TVALVE\)
- Combustion chamber lag \(TGAZ\)

\textbf{Turbine output:}

The turbine mechanical power follows:

\[P_m = BF2 \cdot g_{comb} + AF2\]

where \(g_{comb}\) is the combustion output through lag \(TCD\).

Additional thermodynamic correction terms involve \(TF\) (flame lag) and \(ECR\) (exhaust correction ratio).

\textbf{Parameters from Fortran source (prm array):}

\begin{small}\begin{longtable}[]{@{}lll@{}}
\toprule\noalign{}
Index & Name & Description \\
\midrule\noalign{}
\endhead
\bottomrule\noalign{}
\endlastfoot
1 & \texttt{TDSP} & Speed droop time constant (s) \\
2 & \texttt{STAT} & Steady-state gain \\
3 & \texttt{ZZ} & Governor proportional gain \\
4 & \texttt{XX} & Lead numerator factor \\
5 & \texttt{YY} & Lag denominator factor \\
6 & \texttt{SACC} & Acceleration constant \\
7 & \texttt{MIN} & Minimum fuel valve position (pu) \\
8 & \texttt{MAX} & Maximum fuel valve position (pu) \\
9 & \texttt{TVALVE} & Fuel valve time constant (s) \\
10 & \texttt{TGAZ} & Combustion chamber time constant (s) \\
11 & \texttt{TF} & Flame lag time constant (s) \\
12 & \texttt{ECR} & Exhaust correction ratio \\
13 & \texttt{TCD} & Compressor discharge time constant (s) \\
14 & \texttt{BF2} & Turbine power output slope \\
15 & \texttt{AF2} & Turbine power output intercept (pu) \\
\end{longtable}\end{small}

\textbf{Implementation details:}
- \texttt{nbxvar\ =\ 10}, \texttt{nbzvar\ =\ 4}, \texttt{nbdata\ =\ 15}, \texttt{nbaddpar\ =\ 1}
- Additional: \texttt{prm(16)\ =\ SETPOINT} (initialised from load-flow power)
- Observables: \texttt{g2}, \texttt{x4}, \texttt{g1}

\subsection{Usage Example}\label{model-tor-plant:usage-example}

\textbf{RAMSES Data File}

\begin{Verbatim}
SYNC_MACH  g1  g1
      # FP FQ P Q: share of the power-flow injection (FP, FQ) or initial power (P in MW, Q in Mvar)
      1.  1.  0.  0.
      # SNOM: Nominal apparent power, used as base power in the machine model (MVA)
      800.
      # Pnom: Nominal active power of the turbine, used as base power for the turbine model (MW)
      760.
      # H: Inertia constant (s)
      3.
      # D: Damping coefficient (usually set to zero when the damper windings are modelled) (pu)
      0.
      # IBRATIO: Field current base ratio I_fB^mac/I_fB^exc (see above) (pu)
      0.95
      # TYPE_MOD: Parameter format keyword: RL or XT (case-insensitive)
      XT
      # Xl: Leakage reactance L_l (pu)
      0.15
      # Xd: d-axis synchronous reactance (pu)
      1.1
      # X'd: d-axis transient reactance (must be smaller than Xd) (pu)
      0.25
      # X"d: d-axis subtransient reactance (pu)
      0.2
      # Xq: q-axis synchronous reactance (pu)
      0.7
      # X'q: q-axis transient reactance (pu)
      *
      # X"q: q-axis subtransient reactance (pu)
      0.2
      # m: Saturation coefficient (set to 0 to neglect saturation)
      0.1
      # n: Saturation exponent (ignored when m = 0)
      6.0257
      # Ra: Armature resistance (pu)
      0.
      # T'do: d-axis open-circuit transient time constant (s)
      5.00
      # T"do: d-axis open-circuit subtransient time constant (s)
      0.05
      # T'qo: q-axis open-circuit transient time constant (s)
      *
      # T"qo: q-axis open-circuit subtransient time constant (s)
      0.1
EXC  GENERIC1
      # IFLIM: OEL field current threshold (pu)
      1.8991
      # d: OEL timer band, usually negative: timer runs down at rate 1 below IFLIM+d, holds between IFLIM+d and IFLIM (pu)
      -0.1
      # f: OEL timer charging rate offset above IFLIM
      0.
      # S: OEL timer charging rate per unit of field overcurrent (1/pu)
      1.
      # K1: Not used by the code
      100.
      # K2: OEL gain on (if - IFLIM), negative so that the OEL lowers excitation (pu/pu)
      -1.
      # L1: OEL timer lower limit and initial value, negative; the OEL acts when the timer exceeds 0 (s)
      -11
      # L2: OEL timer upper limit (s)
      10.
      # G: AVR gain (pu/pu)
      70.
      # TA: AVR lead time constant (s)
      10.
      # TB: AVR lag time constant (s)
      20.
      # TE: Exciter time constant (s)
      0.1
      # L3: Minimum field voltage (pu)
      0.
      # L4: Maximum field voltage (pu)
      4.
      # SPEEDIN: PSS input weight: 1 = speed, 0 = active power
      1
      # KPSS: PSS gain, derivative washout KPSS s/(1+sTw); 0 disables PSS (pu s)
      75.
      # Tw: PSS washout time constant (s)
      15.
      # T1: PSS first lead-lag numerator time constant (s)
      0.2
      # T2: PSS first lead-lag denominator time constant (s)
      0.01
      # T3: PSS second lead-lag numerator time constant (s)
      0.2
      # T4: PSS second lead-lag denominator time constant (s)
      0.01
      # DVMIN: PSS output lower limit (pu)
      -0.1
      # DVMAX: PSS output upper limit (pu)
      0.1
TOR  GASTURBM
      # TDSP: rate gain of the load reference driven by SETPOINT (1/s)
      0.05
      # STAT: speed droop (pu)
      0.04
      # ZZ: constant term of lead-lag denominator, static gain is 1/ZZ
      1.0
      # XX: lead time constant of governor lead-lag (s)
      0.50
      # YY: s coefficient of the lead-lag denominator ZZ + YY s; the lag time constant is YY/ZZ (s)
      1.0
      # SACC: acceleration limit reference of acceleration control loop (pu/s)
      1.0
      # MIN: minimum fuel demand (pu)
      0.0
      # MAX: maximum fuel demand (pu)
      1.0
      # TVALVE: fuel valve time constant (s)
      0.05
      # TGAZ: fuel system lag time constant (s)
      0.40
      # TF: second fuel or combustion lag time constant (s)
      0.01
      # ECR: combustion delay, first-order Pade approximation; must be nonzero (s)
      0.01
      # TCD: compressor discharge time constant (s)
      0.20
      # BF2: slope of turbine power versus fuel flow (pu)
      1.0
      # AF2: intercept of turbine power versus fuel flow (pu)
      0.0
;
\end{Verbatim}

\begin{center}\rule{0.5\linewidth}{0.5pt}\end{center}

\section{tor\_govclasm}\label{model-tor-plant:tor_govclasm}

\begin{docsnote}{Note}

Callable since 3.79. Earlier releases compiled this model under no name, so a data file referring to it was rejected.

\end{docsnote}

\subsection{Scientific Description}\label{model-tor-plant:scientific-description-1}

The \texttt{tor\_govclasm} model is a \textbf{classical thermal governor} with a composite control law including deadband, droop, integral control, and combustion dynamics. The name stands for ``governor classical model''.

\textbf{Speed governor with deadband and droop:}

The speed error passes through a deadband \([DBON, DEADB]\) before entering a proportional--integral controller. The integral path has saturation at \(MAXINT\) and gain \(ALPHA\). The reset time \(TRES\) determines the integration rate.

\textbf{Valve and combustion dynamics:}

The valve demand passes through:
1. A governor integrator with time constant \(TSM\), its rate held inside \([VFN,\, V0]\)
2. A first-order lag \(TFV\) (valve actuator)
3. A rate limiter \([PVARMIN,\, PVARMAX]\) (with gain \(GPVAR\)) representing fuel response
4. A combustion/steam chest lag \(KCR\), \(TCR\) producing fuel flow
5. Lower bound: \(MINFUEL\)
6. Transport delay \(TDELAY\) (approximated)
7. Turbine output through lag \(TCH\) with limits \([PMIN,\, PMAX]\)

\textbf{Parameters from Fortran source (prm array):}

\begin{small}\begin{longtable}[]{@{}
  >{\raggedright\arraybackslash}p{(\linewidth - 6\tabcolsep) * \real{0.1465}}
  >{\raggedright\arraybackslash}p{(\linewidth - 6\tabcolsep) * \real{0.2305}}
  >{\raggedright\arraybackslash}p{(\linewidth - 6\tabcolsep) * \real{0.6230}}
@{}}
\toprule\noalign{}
\begin{minipage}[b]{\linewidth}\raggedright
Index
\end{minipage} & \begin{minipage}[b]{\linewidth}\raggedright
Name
\end{minipage} & \begin{minipage}[b]{\linewidth}\raggedright
Description
\end{minipage} \\
\midrule\noalign{}
\endhead
\bottomrule\noalign{}
\endlastfoot
1 & \texttt{STAT} & Governor static characteristic selector \\
2 & \texttt{DEADB} & Deadband half-width (pu) \\
3 & \texttt{DBON} & Deadband onset (pu) \\
4 & \texttt{TSM} & Speed measurement / filter time constant (s) \\
5 & \texttt{V0} & Maximum valve opening rate, upper limit on the governor integrator rate (pu/s) \\
6 & \texttt{VFN} & Minimum valve rate, applied directly as the lower bound, so normally negative (pu/s) \\
7 & \texttt{TFV} & Valve actuator time constant (s) \\
8 & \texttt{VFR} & Fuel flow rate limit (pu/s) \\
9 & \texttt{MAXINT} & Maximum integral output (pu) \\
10 & \texttt{ALPHA} & Integral controller gain \\
11 & \texttt{TRES} & Integral reset time (s) \\
12 & \texttt{GOMP} & Governor proportional gain \\
13 & \texttt{PVARMAX} & Maximum power variation rate (pu) \\
14 & \texttt{PVARMIN} & Minimum power variation rate (pu) \\
15 & \texttt{GPVAR} & Power variation gain \\
16 & \texttt{KCR} & Combustion gain \\
17 & \texttt{TCR} & Combustion time constant (s) \\
18 & \texttt{MINFUEL} & Minimum fuel flow (pu) \\
19 & \texttt{TDELAY} & Combustion transport delay (s) \\
20 & \texttt{PMAX} & Maximum turbine output (pu) \\
21 & \texttt{TCH} & Steam chest / turbine lag (s) \\
22 & \texttt{PMIN} & Minimum turbine output (pu) \\
\end{longtable}\end{small}

\textbf{Implementation details:}
- \texttt{nbxvar\ =\ 6}, \texttt{nbzvar\ =\ 10}, \texttt{nbdata\ =\ 22}, \texttt{nbaddpar\ =\ 1}
- Additional: \texttt{prm(23)\ =\ P0} (initialised from load-flow power)
- Observables: \texttt{x2}, \texttt{b3}, \texttt{b4}

\subsection{Usage Example}\label{model-tor-plant:usage-example-1}

\textbf{RAMSES Data File}

\begin{Verbatim}
SYNC_MACH  g1  g1
      # FP FQ P Q: share of the power-flow injection (FP, FQ) or initial power (P in MW, Q in Mvar)
      1.  1.  0.  0.
      # SNOM: Nominal apparent power, used as base power in the machine model (MVA)
      800.
      # Pnom: Nominal active power of the turbine, used as base power for the turbine model (MW)
      760.
      # H: Inertia constant (s)
      3.
      # D: Damping coefficient (usually set to zero when the damper windings are modelled) (pu)
      0.
      # IBRATIO: Field current base ratio I_fB^mac/I_fB^exc (see above) (pu)
      0.95
      # TYPE_MOD: Parameter format keyword: RL or XT (case-insensitive)
      XT
      # Xl: Leakage reactance L_l (pu)
      0.15
      # Xd: d-axis synchronous reactance (pu)
      1.1
      # X'd: d-axis transient reactance (must be smaller than Xd) (pu)
      0.25
      # X"d: d-axis subtransient reactance (pu)
      0.2
      # Xq: q-axis synchronous reactance (pu)
      0.7
      # X'q: q-axis transient reactance (pu)
      *
      # X"q: q-axis subtransient reactance (pu)
      0.2
      # m: Saturation coefficient (set to 0 to neglect saturation)
      0.1
      # n: Saturation exponent (ignored when m = 0)
      6.0257
      # Ra: Armature resistance (pu)
      0.
      # T'do: d-axis open-circuit transient time constant (s)
      5.00
      # T"do: d-axis open-circuit subtransient time constant (s)
      0.05
      # T'qo: q-axis open-circuit transient time constant (s)
      *
      # T"qo: q-axis open-circuit subtransient time constant (s)
      0.1
EXC  GENERIC1
      # IFLIM: OEL field current threshold (pu)
      1.8991
      # d: OEL timer band, usually negative: timer runs down at rate 1 below IFLIM+d, holds between IFLIM+d and IFLIM (pu)
      -0.1
      # f: OEL timer charging rate offset above IFLIM
      0.
      # S: OEL timer charging rate per unit of field overcurrent (1/pu)
      1.
      # K1: Not used by the code
      100.
      # K2: OEL gain on (if - IFLIM), negative so that the OEL lowers excitation (pu/pu)
      -1.
      # L1: OEL timer lower limit and initial value, negative; the OEL acts when the timer exceeds 0 (s)
      -11
      # L2: OEL timer upper limit (s)
      10.
      # G: AVR gain (pu/pu)
      70.
      # TA: AVR lead time constant (s)
      10.
      # TB: AVR lag time constant (s)
      20.
      # TE: Exciter time constant (s)
      0.1
      # L3: Minimum field voltage (pu)
      0.
      # L4: Maximum field voltage (pu)
      4.
      # SPEEDIN: PSS input weight: 1 = speed, 0 = active power
      1
      # KPSS: PSS gain, derivative washout KPSS s/(1+sTw); 0 disables PSS (pu s)
      75.
      # Tw: PSS washout time constant (s)
      15.
      # T1: PSS first lead-lag numerator time constant (s)
      0.2
      # T2: PSS first lead-lag denominator time constant (s)
      0.01
      # T3: PSS second lead-lag numerator time constant (s)
      0.2
      # T4: PSS second lead-lag denominator time constant (s)
      0.01
      # DVMIN: PSS output lower limit (pu)
      -0.1
      # DVMAX: PSS output upper limit (pu)
      0.1
TOR  GOVCLASM
      # STAT: speed droop (pu)
      0.04
      # DEADB: speed deadband half-width (pu)
      0.003
      # DBON: deadband switch, 1 = deadband applied, 0 = none
      1.
      # TSM: valve servomotor time constant (s)
      0.10
      # V0: maximum valve opening rate (pu/s)
      1.0
      # VFN: minimum valve rate, applied as given so normally negative (pu/s)
      -1.0
      # TFV: valve actuator time constant, read but not used (s)
      0.30
      # VFR: fast-valving rate, never used (fast valving not implemented) (pu/s)
      0.5
      # MAXINT: upper limit of valve position, fuel flow and reheater state (pu)
      1.0
      # ALPHA: HP fraction of turbine power
      1.0
      # TRES: reheater time constant (s)
      5.0
      # GOMP: intercept valve opening gain on valve position
      25.0
      # PVARMAX: upper load bound of sliding-pressure setpoint (pu)
      0.10
      # PVARMIN: lower load bound of sliding-pressure setpoint (pu)
      -0.10
      # GPVAR: pressure setpoint at zero load (pu)
      0.05
      # KCR: boiler pressure controller proportional gain
      1.0
      # TCR: boiler pressure controller integral time (s)
      0.20
      # MINFUEL: minimum fuel flow (pu)
      0.05
      # TDELAY: fuel transport delay, Pade approximation (s)
      0.30
      # PMAX: maximum steam pressure (pu)
      1.05
      # TCH: boiler steam storage time constant (s)
      0.30
      # PMIN: minimum steam pressure (pu)
      0.05
;
\end{Verbatim}

\begin{center}\rule{0.5\linewidth}{0.5pt}\end{center}

\section{tor\_govhydr}\label{model-tor-plant:tor_govhydr}

\begin{docsnote}{Note}

Callable since 3.79. Earlier releases compiled this model under no name, so a data file referring to it was rejected.

\end{docsnote}

\subsection{Scientific Description}\label{model-tor-plant:scientific-description-2}

The \texttt{tor\_govhydr} model is a \textbf{simplified hydraulic governor} with PI control and a non-linear penstock. It differs from \texttt{tor\_hydro\_generic1} in its simplified turbine characteristic and simpler water-column formulation.

\textbf{PI governor:}

\[\dot{x}_2 = K_I \cdot e, \quad x_2 \in [0, 1]\]

\[z_{demand} = K_P \cdot e + x_2\]

where the speed error \(e = P_0/\omega - T_m - \omega \cdot STAT + \ldots\) accounts for the droop and steady-state bias.

\textbf{Penstock (two-lag water column):}

\[T_{CE1} \cdot \frac{dx_3}{dt} = P - x_3(1 + T_{CE2}/T_{CE1})\]

The turbine mechanical power is computed via an empirical quadratic:

\[P_m = 1.35 \cdot z - 0.7 z^2\]

derived from the non-linear head-flow curve approximated at initialisation.

\textbf{Parameters from Fortran source:}

\begin{small}\begin{longtable}[]{@{}
  >{\raggedright\arraybackslash}p{(\linewidth - 6\tabcolsep) * \real{0.1393}}
  >{\raggedright\arraybackslash}p{(\linewidth - 6\tabcolsep) * \real{0.1926}}
  >{\raggedright\arraybackslash}p{(\linewidth - 6\tabcolsep) * \real{0.6681}}
@{}}
\toprule\noalign{}
\begin{minipage}[b]{\linewidth}\raggedright
Index
\end{minipage} & \begin{minipage}[b]{\linewidth}\raggedright
Name
\end{minipage} & \begin{minipage}[b]{\linewidth}\raggedright
Description
\end{minipage} \\
\midrule\noalign{}
\endhead
\bottomrule\noalign{}
\endlastfoot
1 & \texttt{STAT} & Steady-state gain selector \\
2 & \texttt{KP} & Proportional gain \\
3 & \texttt{KI} & Integral gain \\
4 & \texttt{TC} & Governor/servo time constant (s) \\
5 & \texttt{V0} & Maximum gate opening rate, upper limit on gate velocity (pu/s) \\
6 & \texttt{VF} & Maximum gate closing rate, entered as a positive magnitude and applied as \(-V_F\) (pu/s) \\
7 & \texttt{TCE2} & Second penstock time constant (s) \\
8 & \texttt{TCE1} & First penstock time constant (s) \\
\end{longtable}\end{small}

\textbf{Implementation details:}
- \texttt{nbxvar\ =\ 4}, \texttt{nbzvar\ =\ 4}, \texttt{nbdata\ =\ 8}, \texttt{nbaddpar\ =\ 1}
- Additional: \texttt{prm(9)\ =\ P0} (initialised from load-flow power)
- No observables defined (\texttt{nbobs\ =\ 0})

\subsection{Usage Example}\label{model-tor-plant:usage-example-2}

\textbf{RAMSES Data File}

\begin{Verbatim}
SYNC_MACH  g1  g1
      # FP FQ P Q: share of the power-flow injection (FP, FQ) or initial power (P in MW, Q in Mvar)
      1.  1.  0.  0.
      # SNOM: Nominal apparent power, used as base power in the machine model (MVA)
      800.
      # Pnom: Nominal active power of the turbine, used as base power for the turbine model (MW)
      760.
      # H: Inertia constant (s)
      3.
      # D: Damping coefficient (usually set to zero when the damper windings are modelled) (pu)
      0.
      # IBRATIO: Field current base ratio I_fB^mac/I_fB^exc (see above) (pu)
      0.95
      # TYPE_MOD: Parameter format keyword: RL or XT (case-insensitive)
      XT
      # Xl: Leakage reactance L_l (pu)
      0.15
      # Xd: d-axis synchronous reactance (pu)
      1.1
      # X'd: d-axis transient reactance (must be smaller than Xd) (pu)
      0.25
      # X"d: d-axis subtransient reactance (pu)
      0.2
      # Xq: q-axis synchronous reactance (pu)
      0.7
      # X'q: q-axis transient reactance (pu)
      *
      # X"q: q-axis subtransient reactance (pu)
      0.2
      # m: Saturation coefficient (set to 0 to neglect saturation)
      0.1
      # n: Saturation exponent (ignored when m = 0)
      6.0257
      # Ra: Armature resistance (pu)
      0.
      # T'do: d-axis open-circuit transient time constant (s)
      5.00
      # T"do: d-axis open-circuit subtransient time constant (s)
      0.05
      # T'qo: q-axis open-circuit transient time constant (s)
      *
      # T"qo: q-axis open-circuit subtransient time constant (s)
      0.1
EXC  GENERIC1
      # IFLIM: OEL field current threshold (pu)
      1.8991
      # d: OEL timer band, usually negative: timer runs down at rate 1 below IFLIM+d, holds between IFLIM+d and IFLIM (pu)
      -0.1
      # f: OEL timer charging rate offset above IFLIM
      0.
      # S: OEL timer charging rate per unit of field overcurrent (1/pu)
      1.
      # K1: Not used by the code
      100.
      # K2: OEL gain on (if - IFLIM), negative so that the OEL lowers excitation (pu/pu)
      -1.
      # L1: OEL timer lower limit and initial value, negative; the OEL acts when the timer exceeds 0 (s)
      -11
      # L2: OEL timer upper limit (s)
      10.
      # G: AVR gain (pu/pu)
      70.
      # TA: AVR lead time constant (s)
      10.
      # TB: AVR lag time constant (s)
      20.
      # TE: Exciter time constant (s)
      0.1
      # L3: Minimum field voltage (pu)
      0.
      # L4: Maximum field voltage (pu)
      4.
      # SPEEDIN: PSS input weight: 1 = speed, 0 = active power
      1
      # KPSS: PSS gain, derivative washout KPSS s/(1+sTw); 0 disables PSS (pu s)
      75.
      # Tw: PSS washout time constant (s)
      15.
      # T1: PSS first lead-lag numerator time constant (s)
      0.2
      # T2: PSS first lead-lag denominator time constant (s)
      0.01
      # T3: PSS second lead-lag numerator time constant (s)
      0.2
      # T4: PSS second lead-lag denominator time constant (s)
      0.01
      # DVMIN: PSS output lower limit (pu)
      -0.1
      # DVMAX: PSS output upper limit (pu)
      0.1
TOR  GOVHYDR
      # STAT: permanent speed droop (pu)
      0.04
      # KP: PI proportional gain
      2.0
      # KI: PI integral gain (1/s)
      0.5
      # TC: gate servomotor time constant (s)
      0.10
      # V0: maximum gate opening rate (pu/s)
      1.0
      # VF: maximum gate closing rate, positive, applied as minus VF (pu/s)
      1.0
      # TCE2: water column numerator time constant, (1 - s TCE2)/(1 + s TCE1) (s)
      1.0
      # TCE1: water column denominator time constant (s)
      5.0
;
\end{Verbatim}

\textbf{Notes}

\begin{itemize}
\tightlist
\item
  \texttt{TCE1} corresponds roughly to the water-starting time \(T_W\).
\item
  \texttt{TCE2/TCE1} sets the ratio of the two penstock lags.
\end{itemize}

\begin{center}\rule{0.5\linewidth}{0.5pt}\end{center}

\section{tor\_govnuc}\label{model-tor-plant:tor_govnuc}

\begin{docsnote}{Note}

Callable since 3.79. Earlier releases compiled this model under no name, so a data file referring to it was rejected.

\end{docsnote}

\subsection{Scientific Description}\label{model-tor-plant:scientific-description-3}

The \texttt{tor\_govnuc} model is a \textbf{nuclear plant governor} with deadband, droop, integral control, and a non-linear turbine characteristic via \texttt{GOMP} (the gain of the output multiplier). It is structurally similar to \texttt{tor\_govclasm} but tailored for the slower, tightly-regulated dynamics of nuclear steam supply systems.

\textbf{Speed governor:}

The speed error passes through a deadband \([DBON, DEADB]\):

\[e = \text{deadband}(\omega_{ref} - \omega) - STAT \cdot (P - P_0)\]

The error enters an integrator with rate limits \([GRADMIN, GRADMAX]\) (pu/s) filtered through a ramp time \(TGRAD\). The integrator output \(x_I\) passes through a servo \(TSM\).

\textbf{Turbine characteristic with GOMP:}

The turbine output is non-linearly mapped through a gain \(GOMP\):

\[P_m = \min\!\left(\frac{x_I}{GOMP},\, 1\right) \cdot f(x_I)\]

where the limiter prevents over-power and \(f\) is a piecewise-linear characteristic.

The steam-valve rate is held inside \([VFN,\, V0]\), then passes the valve actuator lag \(TFV\), the reset lag \(TRES\) and \(GOMP\) to produce the final torque. \(VFR\) overrides the rate during fast valving, a path the shipped model never takes: the code fixes its selector at 1 and leaves the fast-valving logic unwritten.

\textbf{Parameters from Fortran source:}

\begin{small}\begin{longtable}[]{@{}
  >{\raggedright\arraybackslash}p{(\linewidth - 6\tabcolsep) * \real{0.1410}}
  >{\raggedright\arraybackslash}p{(\linewidth - 6\tabcolsep) * \real{0.2284}}
  >{\raggedright\arraybackslash}p{(\linewidth - 6\tabcolsep) * \real{0.6306}}
@{}}
\toprule\noalign{}
\begin{minipage}[b]{\linewidth}\raggedright
Index
\end{minipage} & \begin{minipage}[b]{\linewidth}\raggedright
Name
\end{minipage} & \begin{minipage}[b]{\linewidth}\raggedright
Description
\end{minipage} \\
\midrule\noalign{}
\endhead
\bottomrule\noalign{}
\endlastfoot
1 & \texttt{STAT} & Governor static characteristic selector \\
2 & \texttt{DEADB} & Deadband half-width (pu) \\
3 & \texttt{DBON} & Deadband onset (pu) \\
4 & \texttt{GRADMAX} & Maximum load gradient (pu/s) \\
5 & \texttt{GRADMIN} & Minimum load gradient / unloading rate (pu/s) \\
6 & \texttt{TGRAD} & Ramp filter time constant (s) \\
7 & \texttt{TSM} & Servo time constant (s) \\
8 & \texttt{V0} & Maximum valve opening rate, upper limit on the steam-valve rate (pu/s) \\
9 & \texttt{VFN} & Minimum valve rate, applied directly as the lower bound, so normally negative (pu/s) \\
10 & \texttt{TFV} & Valve actuator time constant (s) \\
11 & \texttt{VFR} & Valve rate limit (pu/s) \\
12 & \texttt{ALPHA} & Integral gain \\
13 & \texttt{TRES} & Reset time constant (s) \\
14 & \texttt{GOMP} & Output power gain \\
\end{longtable}\end{small}

\textbf{Implementation details:}
- \texttt{nbxvar\ =\ 4}, \texttt{nbzvar\ =\ 8}, \texttt{nbdata\ =\ 14}, \texttt{nbaddpar\ =\ 1}
- Additional: \texttt{prm(15)\ =\ P0} (initialised from load-flow power)
- Observable: \texttt{TM}

\subsection{Usage Example}\label{model-tor-plant:usage-example-3}

\textbf{RAMSES Data File}

\begin{Verbatim}
SYNC_MACH  g1  g1
      # FP FQ P Q: share of the power-flow injection (FP, FQ) or initial power (P in MW, Q in Mvar)
      1.  1.  0.  0.
      # SNOM: Nominal apparent power, used as base power in the machine model (MVA)
      800.
      # Pnom: Nominal active power of the turbine, used as base power for the turbine model (MW)
      760.
      # H: Inertia constant (s)
      3.
      # D: Damping coefficient (usually set to zero when the damper windings are modelled) (pu)
      0.
      # IBRATIO: Field current base ratio I_fB^mac/I_fB^exc (see above) (pu)
      0.95
      # TYPE_MOD: Parameter format keyword: RL or XT (case-insensitive)
      XT
      # Xl: Leakage reactance L_l (pu)
      0.15
      # Xd: d-axis synchronous reactance (pu)
      1.1
      # X'd: d-axis transient reactance (must be smaller than Xd) (pu)
      0.25
      # X"d: d-axis subtransient reactance (pu)
      0.2
      # Xq: q-axis synchronous reactance (pu)
      0.7
      # X'q: q-axis transient reactance (pu)
      *
      # X"q: q-axis subtransient reactance (pu)
      0.2
      # m: Saturation coefficient (set to 0 to neglect saturation)
      0.1
      # n: Saturation exponent (ignored when m = 0)
      6.0257
      # Ra: Armature resistance (pu)
      0.
      # T'do: d-axis open-circuit transient time constant (s)
      5.00
      # T"do: d-axis open-circuit subtransient time constant (s)
      0.05
      # T'qo: q-axis open-circuit transient time constant (s)
      *
      # T"qo: q-axis open-circuit subtransient time constant (s)
      0.1
EXC  GENERIC1
      # IFLIM: OEL field current threshold (pu)
      1.8991
      # d: OEL timer band, usually negative: timer runs down at rate 1 below IFLIM+d, holds between IFLIM+d and IFLIM (pu)
      -0.1
      # f: OEL timer charging rate offset above IFLIM
      0.
      # S: OEL timer charging rate per unit of field overcurrent (1/pu)
      1.
      # K1: Not used by the code
      100.
      # K2: OEL gain on (if - IFLIM), negative so that the OEL lowers excitation (pu/pu)
      -1.
      # L1: OEL timer lower limit and initial value, negative; the OEL acts when the timer exceeds 0 (s)
      -11
      # L2: OEL timer upper limit (s)
      10.
      # G: AVR gain (pu/pu)
      70.
      # TA: AVR lead time constant (s)
      10.
      # TB: AVR lag time constant (s)
      20.
      # TE: Exciter time constant (s)
      0.1
      # L3: Minimum field voltage (pu)
      0.
      # L4: Maximum field voltage (pu)
      4.
      # SPEEDIN: PSS input weight: 1 = speed, 0 = active power
      1
      # KPSS: PSS gain, derivative washout KPSS s/(1+sTw); 0 disables PSS (pu s)
      75.
      # Tw: PSS washout time constant (s)
      15.
      # T1: PSS first lead-lag numerator time constant (s)
      0.2
      # T2: PSS first lead-lag denominator time constant (s)
      0.01
      # T3: PSS second lead-lag numerator time constant (s)
      0.2
      # T4: PSS second lead-lag denominator time constant (s)
      0.01
      # DVMIN: PSS output lower limit (pu)
      -0.1
      # DVMAX: PSS output upper limit (pu)
      0.1
TOR  GOVNUC
      # STAT: speed droop (pu)
      0.04
      # DEADB: speed deadband half-width (pu)
      0.005
      # DBON: deadband switch, 1 = deadband applied, 0 = none
      1.
      # GRADMAX: upper clamp on load reference error, ramp rate is GRADMAX/TGRAD (pu)
      0.02
      # GRADMIN: lower clamp magnitude, positive, applied as minus GRADMIN (pu)
      0.02
      # TGRAD: load reference ramp time constant (s)
      10.0
      # TSM: valve servomotor time constant (s)
      5.0
      # V0: maximum valve opening rate (pu/s)
      1.0
      # VFN: minimum valve rate, applied as given so normally negative (pu/s)
      -1.0
      # TFV: valve actuator time constant, read but not used (s)
      0.50
      # VFR: fast-valving rate, never used (fast valving not implemented) (pu/s)
      0.10
      # ALPHA: HP fraction of turbine power
      1.0
      # TRES: reheater (moisture separator) time constant (s)
      20.0
      # GOMP: intercept valve opening gain on valve position
      1.0
;
\end{Verbatim}

\textbf{Notes}

\begin{itemize}
\tightlist
\item
  Nuclear plants typically have slow load-following characteristics: \texttt{GRADMAX} is typically 0.005--0.02 pu/s.
\item
  The \texttt{DEADB} is important to prevent hunting around the setpoint.
\item
  \texttt{GOMP} scales the turbine output; set to 1.0 for standard normalisation.
\end{itemize}
